\documentclass[../../main.tex]{subfiles}

\begin{document}

\chapter{Clases de la API}


%===========================================================
%===========================================================
%                       Clase System
%===========================================================
%===========================================================

\section{Clase System}

La clase \texttt{System} constituye el punto de acceso principal a las funciones globales del controlador Suruga Seiki. A través de esta clase es posible consultar las capacidades del sistema, verificar el estado de la comunicación con el controlador, obtener información sobre la versión del software y gestionar operaciones generales relacionadas con la conexión y el diagnóstico del equipo.

La Tabla~\ref{tab:system_members} resume los principales miembros públicos de la clase, agrupados en propiedades, métodos, campos y tipos anidados. Para cada elemento se indica su propósito, el tipo de dato asociado y el nivel de acceso proporcionado por la API.

%================================================================
%================================================================
%================================================================

\footnotesize

\begin{longtblr}[
    caption={Miembros de la clase \texttt{System}},
    label={tab:system_members},
]{
    width=\textwidth,
    colspec={
    |Q[c,wd=0.28\textwidth]|
    X|
    Q[c,wd=0.12\textwidth]|
    Q[c,wd=0.10\textwidth]|
    },
    hlines,
}

\SetRow{bg=gray!50}
\textbf{Miembro} &
\textbf{Descripción} &
\textbf{Tipo .NET} &
\textbf{Acceso} \\

%===========================================================
% Propiedades
%===========================================================
\SetRow{bg=gray9}
\SetCell[c=4]{c}
Propiedades \\

MaxAxisStageNumber &
Número máximo de ejes de posicionamiento soportados por el sistema. &
UInt16 &
get/set \\

MaxAnalogInputChNumber &
Número máximo de canales de entrada analógica disponibles. &
UInt16 &
get/set \\

MaxAlignmentNumber &
Número máximo de instancias de alineación soportadas por el sistema. &
UInt16 &
get/set \\

MaxAngleAdjustmentNumber &
Número máximo de módulos de ajuste angular soportados. &
UInt16 &
get/set \\

MaxDInChNumber &
Número máximo de canales de entrada digital disponibles. &
UInt16 &
get/set \\

MaxDOutChNumber &
Número máximo de canales de salida digital disponibles. &
UInt16 &
get/set \\

MaxAnalogOutputChNumber &
Número máximo de canales de salida analógica disponibles. &
UInt16 &
get/set \\

MaxPowerMeterChNumber &
Número máximo de canales del medidor de potencia soportados. &
UInt16 &
get/set \\

Instance &
Instancia única (Singleton) de la clase \texttt{System}. &
System &
get \\

Connected &
Indica si existe una conexión activa con el controlador. &
Boolean &
get \\

DllVersion &
Versión de la biblioteca dinámica (DLL) utilizada por la API. &
String &
get \\

SystemVersion &
Versión del firmware o software del sistema de control. &
String &
get \\

%===========================================================
% Métodos
%===========================================================
\SetRow{bg=gray9}
\SetCell[c=4]{c}
Métodos \\

IsError() &
Determina si el sistema presenta algún estado de error. &
Boolean &
-- \\

GetCurrentErrorCode() &
Obtiene el código del error actual del sistema. &
ErrorCode &
-- \\

IsEmergencyAsserted() &
Indica si la condición de parada de emergencia está activada. &
Boolean &
-- \\

SetAddress(String addr) &
Establece la dirección del controlador utilizada para la comunicación. &
Void &
-- \\

DisconnectAds() &
Finaliza la conexión ADS con el controlador. &
Void &
-- \\

%===========================================================
% Campos
%===========================================================
\SetRow{bg=gray9}
\SetCell[c=4]{c}
Campos \\

MaxProfileNumber &
Número máximo de perfiles soportados por el sistema. &
UInt16 &
static \\

ProfileDataPacketMaxDataNumber &
Número máximo de datos que puede contener un paquete de datos de perfil. &
UInt32 &
static \\

%===========================================================
% Tipos anidados / Enums
%===========================================================
\SetRow{bg=gray9}
\SetCell[c=4]{c}
Tipos anidados / Enums \\

ErrorCode &
Enumeración que define los posibles códigos de error devueltos por la clase \texttt{System}. &
Enum &
-- \\

\end{longtblr}

\normalsize

%================================================================
%================================================================
%================================================================



%===========================================================
%===========================================================
%                   Clase AxisComponents
%===========================================================
%===========================================================

\section{Clase AxisComponents}

%================================================================
%================================================================
%================================================================

\footnotesize

\begin{longtblr}[
    caption={Miembros de la clase \texttt{AxisComponents}},
    label={tab:axiscomponents_members},
]{
    width=\textwidth,
    colspec={
    |Q[c,wd=0.28\textwidth]|
    X|
    Q[c,wd=0.12\textwidth]|
    Q[c,wd=0.10\textwidth]|
    },
    hlines,
}

\SetRow{bg=gray!50}
\textbf{Miembro} &
\textbf{Descripción} &
\textbf{Tipo .NET} &
\textbf{Acceso} \\

%===========================================================
% Constructores
%===========================================================
\SetRow{bg=gray9}
\SetCell[c=4]{c}
Constructores \\

AxisComponents(UInt16 axisNumber) &
Inicializa una nueva instancia de la clase asociada al eje especificado. &
Constructor &
-- \\

%===========================================================
% Métodos
%===========================================================
\SetRow{bg=gray9}
\SetCell[c=4]{c}
Métodos \\

MoveAbsolute(Double position) &
Desplaza el eje hasta una posición absoluta. &
ErrorCode &
-- \\

MoveRelative(Double distance) &
Desplaza el eje una distancia relativa respecto a su posición actual. &
ErrorCode &
-- \\

ReturnOrigin() &
Ejecuta el procedimiento de retorno al origen del eje. &
ErrorCode &
-- \\

Stop() &
Detiene inmediatamente el movimiento del eje. &
ErrorCode &
-- \\

TurnOnServo() &
Activa el servomotor del eje. &
ErrorCode &
-- \\

TurnOffServo() &
Desactiva el servomotor del eje. &
ErrorCode &
-- \\

ResetError() &
Restablece el estado de error del eje. &
ErrorCode &
-- \\

GetTargetPosition() &
Obtiene la posición objetivo configurada para el eje. &
Double &
-- \\

GetActualPosition() &
Obtiene la posición actual del eje. &
Double &
-- \\

IsServoOn() &
Indica si el servomotor se encuentra activado. &
Boolean &
-- \\

IsServoAlarm() &
Indica si el servomotor presenta una condición de alarma. &
Boolean &
-- \\

IsMoving() &
Indica si el eje se encuentra en movimiento. &
Boolean &
-- \\

IsOriginSignalOn() &
Indica si la señal del sensor de origen está activa. &
Boolean &
-- \\

IsCwLimitSignalOn() &
Indica si el interruptor de límite en sentido horario (CW) está activo. &
Boolean &
-- \\

IsCcwLimitSignalOn() &
Indica si el interruptor de límite en sentido antihorario (CCW) está activo. &
Boolean &
-- \\

IsCwTorqueLimitOn() &
Indica si se ha alcanzado el límite de par en sentido horario. &
Boolean &
-- \\

IsCcwTorqueLimitOn() &
Indica si se ha alcanzado el límite de par en sentido antihorario. &
Boolean &
-- \\

IsOriginReturnSucceeded() &
Indica si el procedimiento de retorno al origen finalizó correctamente. &
Boolean &
-- \\

SetAccelRate(Double rate) &
Configura la tasa de aceleración del eje. &
ErrorCode &
-- \\

SetDecelRate(Double rate) &
Configura la tasa de desaceleración del eje. &
ErrorCode &
-- \\

SetSineMotion(Boolean sineMotion) &
Habilita o deshabilita el modo de movimiento senoidal. &
ErrorCode &
-- \\

SetFulcrumSpan(Double span) &
Configura la distancia del brazo o fulcro utilizada durante el movimiento. &
ErrorCode &
-- \\

SetEncoderStepLength(Double length) &
Configura la longitud correspondiente a un paso del codificador. &
ErrorCode &
-- \\

SetMaxSpeed(Double speed) &
Establece la velocidad máxima permitida para el eje. &
ErrorCode &
-- \\

SetCWTorqueLimit(Double cWlimit) &
Configura el límite de par en sentido horario (CW). &
ErrorCode &
-- \\

SetCCWTorqueLimit(Double cCWlimit) &
Configura el límite de par en sentido antihorario (CCW). &
ErrorCode &
-- \\

GetErrorCode() &
Obtiene el código de error actual del eje. &
ErrorCode &
-- \\

GetStatus() &
Obtiene el estado actual del eje. &
Status &
-- \\

%===========================================================
% Tipos anidados / Enums
%===========================================================
\SetRow{bg=gray9}
\SetCell[c=4]{c}
Tipos anidados / Enums \\

ErrorCode &
Enumeración que define los posibles códigos de error devueltos por la clase \texttt{AxisComponents}. &
Enum &
-- \\

Status &
Enumeración que representa los distintos estados de funcionamiento del eje. &
Enum &
-- \\

\end{longtblr}

\normalsize

%================================================================
%================================================================
%================================================================


%===========================================================
%===========================================================
%                       Clase Axis2D
%===========================================================
%===========================================================

\section{Clase Axis2D}

%================================================================
%================================================================
%================================================================

\footnotesize

\begin{longtblr}[
    caption={Miembros de la clase \texttt{Axis2D}},
    label={tab:axis2d_members},
]{
    width=\textwidth,
    colspec={
    |Q[c,wd=0.28\textwidth]|
    X|
    Q[c,wd=0.12\textwidth]|
    Q[c,wd=0.10\textwidth]|
    },
    hlines,
}

\SetRow{bg=gray!50}
\textbf{Miembro} &
\textbf{Descripción} &
\textbf{Tipo .NET} &
\textbf{Acceso} \\

%===========================================================
% Constructores
%===========================================================
\SetRow{bg=gray9}
\SetCell[c=4]{c}
Constructores \\

Axis2D(UInt16 mainStageNumber, UInt16 subStageNumber) &
Inicializa una nueva instancia de la clase asociada a un sistema de posicionamiento bidimensional compuesto por dos ejes. &
Constructor &
-- \\

%===========================================================
% Métodos
%===========================================================
\SetRow{bg=gray9}
\SetCell[c=4]{c}
Métodos \\

SetAxisNumber(UInt16 mainStageNumber, UInt16 subStageNumber) &
Asigna los identificadores de los ejes principal y secundario utilizados por el sistema bidimensional. &
ErrorCode &
-- \\

GetAxisNumber() &
Obtiene los identificadores de los ejes asociados al sistema bidimensional. &
UInt16[] &
-- \\

MoveAbsolute(Point targetPosition) &
Desplaza ambos ejes hasta una posición absoluta especificada mediante una estructura de coordenadas. &
ErrorCode &
-- \\

MoveAbsolute(Double mainTargetPosition, Double subTargetPosition) &
Desplaza ambos ejes hasta las posiciones absolutas indicadas para cada eje. &
ErrorCode &
-- \\

MoveRelative(Point targetDistance) &
Realiza un desplazamiento relativo de ambos ejes utilizando una estructura de coordenadas. &
ErrorCode &
-- \\

MoveRelative(Double mainDistance, Double subDistance) &
Realiza un desplazamiento relativo independiente para cada eje. &
ErrorCode &
-- \\

MoveAbsoluteAngle(Double mainTargetPosition, Double subAngle) &
Desplaza el eje principal a una posición absoluta mientras posiciona el eje secundario según un ángulo especificado. &
ErrorCode &
-- \\

MoveRelativeAngle(Double mainDistance, Double subAngle) &
Realiza un desplazamiento relativo del eje principal y posiciona el eje secundario mediante un ángulo especificado. &
ErrorCode &
-- \\

Stop() &
Detiene el movimiento de ambos ejes. &
ErrorCode &
-- \\

GetTargetPosition() &
Obtiene la posición objetivo del sistema bidimensional. &
Point &
-- \\

GetActualPosition() &
Obtiene la posición actual del sistema bidimensional. &
Point &
-- \\

IsMoving() &
Indica si alguno de los ejes se encuentra en movimiento. &
Boolean &
-- \\

SetSpeed(Double speed) &
Configura la velocidad de desplazamiento del sistema bidimensional. &
ErrorCode &
-- \\

SetAccelRate(Double rate) &
Configura la tasa de aceleración del movimiento. &
ErrorCode &
-- \\

SetDecelRate(Double rate) &
Configura la tasa de desaceleración del movimiento. &
ErrorCode &
-- \\

GetErrorCode() &
Obtiene el código de error actual del sistema bidimensional. &
ErrorCode &
-- \\

GetStatus() &
Obtiene el estado actual del sistema bidimensional. &
Status &
-- \\

%===========================================================
% Tipos anidados / Enums
%===========================================================
\SetRow{bg=gray9}
\SetCell[c=4]{c}
Tipos anidados / Enums \\

Point &
Estructura que representa una posición bidimensional mediante las coordenadas de los ejes principal y secundario. &
Struct &
-- \\

\end{longtblr}

\normalsize

%================================================================
%================================================================
%================================================================


%===========================================================
%===========================================================
%                       Clase Axis3D
%===========================================================
%===========================================================

\section{Clase Axis3D}

%================================================================
%================================================================
%================================================================

\footnotesize

\begin{longtblr}[
    caption={Miembros de la clase \texttt{Axis3D}},
    label={tab:axis3d_members},
]{
    width=\textwidth,
    colspec={
        |Q[l,wd=0.28\textwidth]|
        X|
        Q[c,wd=0.12\textwidth]|
        Q[c,wd=0.10\textwidth]|
    },
    hlines,
}

\SetRow{bg=gray!50}
\textbf{Miembro} &
\textbf{Descripción} &
\textbf{Tipo .NET} &
\textbf{Acceso} \\

%===========================================================
% Constructores
%===========================================================
\SetRow{bg=gray9}
\SetCell[c=4]{c}
Constructores \\

Axis3D(UInt16 mainStageNumberX, UInt16 mainStageNumberY, UInt16 mainStageNumberZ) &
Inicializa una nueva instancia de la clase asociada a un sistema de posicionamiento tridimensional compuesto por tres ejes. &
Constructor &
-- \\

%===========================================================
% Métodos
%===========================================================
\SetRow{bg=gray9}
\SetCell[c=4]{c}
Métodos \\

SetAxisNumber(UInt16 mainStageNumberX, UInt16 mainStageNumberY, UInt16 mainStageNumberZ) &
Asigna los identificadores de los ejes utilizados por el sistema tridimensional. &
ErrorCode &
-- \\

GetAxisNumber() &
Obtiene los identificadores de los ejes asociados al sistema tridimensional. &
UInt16[] &
-- \\

MoveAbsolute(Point3D targetPosition) &
Desplaza el sistema tridimensional hasta una posición absoluta especificada mediante una estructura de coordenadas. &
ErrorCode &
-- \\

MoveAbsolute(Double targetPositionX, Double targetPositionY, Double targetPositionZ) &
Desplaza el sistema tridimensional hasta las posiciones absolutas indicadas para cada eje. &
ErrorCode &
-- \\

MoveAbsolute(Point3D targetPosition, UInt16 dOutId, Int16 dOutOnDelayTime, Int16 dOutOffDelayTime) &
Desplaza el sistema tridimensional hasta una posición absoluta activando una salida digital durante el movimiento según los tiempos especificados. &
ErrorCode &
-- \\

MoveAbsolute(Double targetPositionX, Double targetPositionY, Double targetPositionZ, UInt16 dOutId, Int16 dOutOnDelayTime, Int16 dOutOffDelayTime) &
Desplaza el sistema tridimensional hasta las posiciones absolutas indicadas para cada eje activando una salida digital durante el movimiento. &
ErrorCode &
-- \\

MoveRelative(Point3D targetDistance) &
Realiza un desplazamiento relativo del sistema tridimensional utilizando una estructura de coordenadas. &
ErrorCode &
-- \\

MoveRelative(Double targetDistanceX, Double targetDistanceY, Double targetDistanceZ) &
Realiza un desplazamiento relativo independiente para cada eje. &
ErrorCode &
-- \\

MoveRelative(Point3D targetDistance, UInt16 dOutId, Int16 dOutOnDelayTime, Int16 dOutOffDelayTime) &
Realiza un desplazamiento relativo activando una salida digital durante el movimiento según los tiempos especificados. &
ErrorCode &
-- \\

MoveRelative(Double targetDistanceX, Double targetDistanceY, Double targetDistanceZ, UInt16 dOutId, Int16 dOutOnDelayTime, Int16 dOutOffDelayTime) &
Realiza un desplazamiento relativo independiente para cada eje activando una salida digital durante el movimiento. &
ErrorCode &
-- \\

Stop() &
Detiene el movimiento de los tres ejes. &
ErrorCode &
-- \\

GetTargetPosition() &
Obtiene la posición objetivo del sistema tridimensional. &
Point3D &
-- \\

GetActualPosition() &
Obtiene la posición actual del sistema tridimensional. &
Point3D &
-- \\

IsMoving() &
Indica si alguno de los ejes se encuentra en movimiento. &
Boolean &
-- \\

SetSpeed(Double speed) &
Configura la velocidad de desplazamiento del sistema tridimensional. &
ErrorCode &
-- \\

SetAccelRate(Double accelRate) &
Configura la tasa de aceleración del movimiento. &
ErrorCode &
-- \\

SetDecelRate(Double decelRate) &
Configura la tasa de desaceleración del movimiento. &
ErrorCode &
-- \\

SetStageNumber(UInt16 mainStageNumberX, UInt16 mainStageNumberY, UInt16 mainStageNumberZ) &
Configura los identificadores de los ejes principales utilizados por el sistema. &
ErrorCode &
-- \\

SetSubStageNumber(UInt16 subStageNumberX, UInt16 subStageNumberY, UInt16 subStageNumberZ) &
Configura los identificadores de los ejes secundarios utilizados por el sistema. &
ErrorCode &
-- \\

SetRotationCenter(Double rotationOffsetX, Double rotationOffsetY, Double rotationOffsetZ) &
Define el centro de rotación del sistema mediante los desplazamientos especificados para los ejes X, Y y Z. &
ErrorCode &
-- \\

GetErrorCode() &
Obtiene el código de error actual del sistema tridimensional. &
ErrorCode &
-- \\

GetStatus() &
Obtiene el estado actual del sistema tridimensional. &
Status &
-- \\

%===========================================================
% Tipos anidados / Estructuras
%===========================================================
\SetRow{bg=gray9}
\SetCell[c=4]{c}
Tipos anidados / Estructuras \\

Point3D &
Estructura que representa una posición tridimensional mediante las coordenadas X, Y y Z. &
Struct &
-- \\

RotationCenter &
Estructura que define los ejes y desplazamientos utilizados como centro de rotación del sistema. &
Struct &
-- \\

\end{longtblr}

\normalsize

%================================================================
%================================================================
%================================================================

%===========================================================
%===========================================================
%                     Clase Alignment
%===========================================================
%===========================================================

\section{Clase Alignment}

%================================================================
%================================================================
%================================================================

\footnotesize

\begin{longtblr}[
    caption={Miembros de la clase \texttt{Alignment}},
    label={tab:alignment_members},
]{
    width=\textwidth,
    colspec={
        |Q[l,wd=0.28\textwidth]|
        X|
        Q[c,wd=0.12\textwidth]|
        Q[c,wd=0.10\textwidth]|
    },
    hlines,
}

\SetRow{bg=gray!50}
\textbf{Miembro} &
\textbf{Descripción} &
\textbf{Tipo .NET} &
\textbf{Acceso} \\

%===========================================================
% Constructores
%===========================================================
\SetRow{bg=gray9}
\SetCell[c=4]{c}
Constructores \\

Alignment() &
Inicializa una nueva instancia de la clase \texttt{Alignment}. &
Constructor &
-- \\

Alignment(UInt16 id) &
Inicializa una nueva instancia de la clase asociada al identificador de alineación especificado. &
Constructor &
-- \\

%===========================================================
% Métodos
%===========================================================
\SetRow{bg=gray9}
\SetCell[c=4]{c}
Métodos \\

StartSingle() &
Inicia el procedimiento de alineación de tipo \emph{Single}. &
ErrorCode &
-- \\

StartFlat() &
Inicia el procedimiento de alineación de tipo \emph{Flat}. &
ErrorCode &
-- \\

StartFocus() &
Inicia el procedimiento de alineación de tipo \emph{Focus}. &
ErrorCode &
-- \\

StartTwoChRotation() &
Inicia el procedimiento de alineación por rotación de dos canales. &
ErrorCode &
-- \\

StartMMFRotation() &
Inicia el procedimiento de alineación por rotación para fibra multimodo (MMF). &
ErrorCode &
-- \\

StartLoopBack() &
Inicia el procedimiento de alineación de tipo \emph{LoopBack}. &
ErrorCode &
-- \\

StartCommon() &
Inicia un procedimiento de alineación utilizando parámetros comunes. &
ErrorCode &
-- \\

Stop() &
Detiene el procedimiento de alineación en ejecución. &
ErrorCode &
-- \\

GetErrorCode() &
Obtiene el código de error actual del proceso de alineación. &
ErrorCode &
-- \\

GetStatus() &
Obtiene el estado actual del proceso de alineación. &
Status &
-- \\

GetErrorAxisID() &
Obtiene el identificador del eje responsable del error actual. &
UInt32 &
-- \\

GetAligningStatus() &
Obtiene el estado detallado del procedimiento de alineación. &
StatusCode &
-- \\

GetProfilePacketSumIndex (ProfileDataType profileDataType) &
Obtiene el número total de paquetes de datos disponibles para el perfil especificado. &
UInt32 &
-- \\

RequestProfileData(ProfileDataType profileDataType, UInt32 index) &
Solicita un paquete de datos correspondiente al perfil y al índice especificados. &
ProfileData &
-- \\

RequestProfileData(ProfileDataType profileDataType, UInt32 index, Boolean isRotationCenterShift) &
Solicita un paquete de datos del perfil considerando el desplazamiento del centro de rotación. &
ProfileData &
-- \\

GetVoltage(UInt32 alignmentAInCh) &
Obtiene el valor de tensión medido en el canal de entrada analógica especificado. &
Double &
-- \\

GetPower(UInt32 alignmentPMCh) &
Obtiene la potencia óptica medida por el canal del medidor de potencia especificado. &
Double &
-- \\

GetMeasurementWaveLength(UInt32 alignmentPMCh) &
Obtiene la longitud de onda configurada para el canal del medidor de potencia especificado. &
Double &
-- \\

SetSingle(ISingleParameter parameter) &
Configura los parámetros del procedimiento de alineación \emph{Single}. &
ErrorCode &
-- \\

SetSingle(ISingleParameter parameter, RotationCenter rotationCenter) &
Configura los parámetros del procedimiento \emph{Single} utilizando un centro de rotación definido. &
ErrorCode &
-- \\

SetFlat(IFlatParameter parameter) &
Configura los parámetros del procedimiento de alineación \emph{Flat}. &
ErrorCode &
-- \\

SetFlat(IFlatParameter parameter, RotationCenter rotationCenter) &
Configura los parámetros del procedimiento \emph{Flat} utilizando un centro de rotación definido. &
ErrorCode &
-- \\

SetFocus(IFocusParameter parameter) &
Configura los parámetros del procedimiento de alineación \emph{Focus}. &
ErrorCode &
-- \\

SetFocus(IFocusParameter parameter, RotationCenter rotationCenter) &
Configura los parámetros del procedimiento \emph{Focus} utilizando un centro de rotación definido. &
ErrorCode &
-- \\

SetTwoChRotation(ITwoChRotationParameter parameter) &
Configura los parámetros del procedimiento de alineación por rotación de dos canales. &
ErrorCode &
-- \\

SetTwoChRotation(ITwoChRotationParameter parameter, RotationCenter rotationCenter) &
Configura los parámetros del procedimiento de rotación de dos canales utilizando un centro de rotación definido. &
ErrorCode &
-- \\

SetMMFRotation(IMMFRotationParameter parameter) &
Configura los parámetros del procedimiento de alineación por rotación para fibra multimodo. &
ErrorCode &
-- \\

SetMMFRotation(IMMFRotationParameter parameter, RotationCenter rotationCenter) &
Configura los parámetros del procedimiento de rotación para fibra multimodo utilizando un centro de rotación definido. &
ErrorCode &
-- \\

SetLoopBack(ILoopBackParameter parameter) &
Configura los parámetros del procedimiento de alineación \emph{LoopBack}. &
ErrorCode &
-- \\

SetLoopBack(ILoopBackParameter parameter, RotationCenter rotationCenter) &
Configura los parámetros del procedimiento \emph{LoopBack} utilizando un centro de rotación definido. &
ErrorCode &
-- \\

SetCommon(CommonParameter parameter) &
Configura los parámetros comunes del procedimiento de alineación. &
ErrorCode &
-- \\

SetCommon(CommonParameter parameter, RotationCenter rotationCenter) &
Configura los parámetros comunes utilizando un centro de rotación definido. &
ErrorCode &
-- \\

SetMeasurementWaveLength(UInt32 alignmentPMCh, Double waveLength) &
Configura la longitud de onda de medición para el canal del medidor de potencia especificado. &
ErrorCode &
-- \\

%===========================================================
% Tipos anidados
%===========================================================
\SetRow{bg=gray9}
\SetCell[c=4]{c}
Tipos anidados \\

AlignmentType &
Enumeración que define los distintos tipos de procedimientos de alineación soportados. &
Enum &
-- \\

ErrorCode &
Enumeración que define los posibles códigos de error del proceso de alineación. &
Enum &
-- \\

Status &
Enumeración que representa los estados posibles del procedimiento de alineación. &
Enum &
-- \\

AligningStatusCode &
Enumeración que describe el estado interno del algoritmo de alineación. &
Enum &
-- \\

ProfileDataType &
Enumeración que identifica los diferentes tipos de perfiles generados durante la alineación. &
Enum &
-- \\

ZMode &
Enumeración que define el modo de búsqueda utilizado sobre el eje Z. &
Enum &
-- \\

ISharingParameter &
Interfaz que define los parámetros compartidos entre distintos procedimientos de alineación. &
Interface &
-- \\

ISingleParameter &
Interfaz que define los parámetros del procedimiento \emph{Single}. &
Interface &
-- \\

SingleParameter &
Clase que implementa los parámetros del procedimiento \emph{Single}. &
Class &
-- \\

IFlatParameter &
Interfaz que define los parámetros del procedimiento \emph{Flat}. &
Interface &
-- \\

FlatParameter &
Clase que implementa los parámetros del procedimiento \emph{Flat}. &
Class &
-- \\

IFocusParameter &
Interfaz que define los parámetros del procedimiento \emph{Focus}. &
Interface &
-- \\

FocusParameter &
Clase que implementa los parámetros del procedimiento \emph{Focus}. &
Class &
-- \\

ITwoChRotationParameter &
Interfaz que define los parámetros del procedimiento de rotación de dos canales. &
Interface &
-- \\

TwoChRotationParameter &
Clase que implementa los parámetros del procedimiento de rotación de dos canales. &
Class &
-- \\

IMMFRotationParameter &
Interfaz que define los parámetros del procedimiento de rotación para fibra multimodo. &
Interface &
-- \\

MMFRotationParameter &
Clase que implementa los parámetros del procedimiento de rotación para fibra multimodo. &
Class &
-- \\

ILoopBackParameter &
Interfaz que define los parámetros del procedimiento \emph{LoopBack}. &
Interface &
-- \\

LoopBackParameter &
Clase que implementa los parámetros del procedimiento \emph{LoopBack}. &
Class &
-- \\

ICommonParameter &
Interfaz que define los parámetros comunes de los procedimientos de alineación. &
Interface &
-- \\

CommonParameter &
Clase que implementa los parámetros comunes de alineación. &
Class &
-- \\

ProfileData &
Clase que almacena los datos de perfil obtenidos durante un procedimiento de alineación. &
Class &
-- \\

\end{longtblr}

\normalsize

%================================================================
%================================================================
%================================================================


%===========================================================
%===========================================================
%                     Clase Profile
%===========================================================
%===========================================================

\section{Clase Profile}

%================================================================
%================================================================
%================================================================

\footnotesize

\begin{longtblr}[
    caption={Miembros de la clase \texttt{Profile}},
    label={tab:profile_members},
]{
    width=\textwidth,
    colspec={
        |Q[l,wd=0.28\textwidth]|
        X|
        Q[c,wd=0.12\textwidth]|
        Q[c,wd=0.10\textwidth]|
    },
    hlines,
}

\SetRow{bg=gray!50}
\textbf{Miembro} &
\textbf{Descripción} &
\textbf{Tipo .NET} &
\textbf{Acceso} \\

%===========================================================
% Constructores
%===========================================================
\SetRow{bg=gray9}
\SetCell[c=4]{c}
Constructores \\

Profile() &
Inicializa una nueva instancia de la clase \texttt{Profile}. &
Constructor &
-- \\

%===========================================================
% Métodos
%===========================================================
\SetRow{bg=gray9}
\SetCell[c=4]{c}
Métodos \\

Start() &
Inicia el proceso de adquisición de datos de perfil. &
ErrorCode &
-- \\

Stop() &
Detiene el proceso de adquisición de datos de perfil. &
ErrorCode &
-- \\

GetProfileStatus() &
Obtiene el estado actual del proceso de adquisición del perfil. &
Status &
-- \\

GetProfilePacketSumIndex() &
Obtiene el número total de paquetes de datos disponibles para el perfil adquirido. &
UInt32 &
-- \\

RequestProfileData(UInt32 index) &
Solicita un paquete de datos de perfil correspondiente al índice especificado. &
ProfileData &
-- \\

RequestProfileData(UInt32 index, Boolean isRotationCenterShift) &
Solicita un paquete de datos de perfil considerando el desplazamiento del centro de rotación. &
ProfileData &
-- \\

SetProfile(ProfileParameter parameter) &
Configura los parámetros utilizados para la adquisición del perfil. &
ErrorCode &
-- \\

SetProfile(ProfileParameter parameter, RotationCenter rotationCenter) &
Configura los parámetros de adquisición del perfil utilizando un centro de rotación definido. &
ErrorCode &
-- \\

%===========================================================
% Tipos anidados
%===========================================================
\SetRow{bg=gray9}
\SetCell[c=4]{c}
Tipos anidados \\

ErrorCode &
Enumeración que define los posibles códigos de error del proceso de adquisición del perfil. &
Enum &
-- \\

Status &
Enumeración que representa los estados posibles del proceso de adquisición del perfil. &
Enum &
-- \\

ProfileParameter &
Clase que define los parámetros utilizados para configurar la adquisición de un perfil. &
Class &
-- \\

ProfileData &
Clase que almacena los datos obtenidos durante la adquisición de un perfil. &
Class &
-- \\

\end{longtblr}

\normalsize

%================================================================
%================================================================
%================================================================


%===========================================================
%===========================================================
%                 Clase Angle Adjustment
%===========================================================
%===========================================================

\section{Clase Angle Asjustment}

%================================================================
%================================================================
%================================================================

\footnotesize

\begin{longtblr}[
    caption={Miembros de la clase \texttt{AngleAdjustment}},
    label={tab:angleadjustment_members},
]{
    width=\textwidth,
    colspec={
        |Q[l,wd=0.28\textwidth]|
        X|
        Q[c,wd=0.12\textwidth]|
        Q[c,wd=0.10\textwidth]|
    },
    hlines,
}

\SetRow{bg=gray!50}
\textbf{Miembro} &
\textbf{Descripción} &
\textbf{Tipo .NET} &
\textbf{Acceso} \\

%===========================================================
% Constructores
%===========================================================
\SetRow{bg=gray9}
\SetCell[c=4]{c}
Constructores \\

AngleAdjustment(UInt16 id) &
Inicializa una nueva instancia de la clase asociada al identificador del módulo de ajuste angular especificado. &
Constructor &
-- \\

%===========================================================
% Métodos
%===========================================================
\SetRow{bg=gray9}
\SetCell[c=4]{c}
Métodos \\

Start() &
Inicia el procedimiento de ajuste angular. &
ErrorCode &
-- \\

Stop() &
Detiene el procedimiento de ajuste angular. &
ErrorCode &
-- \\

GetStatus() &
Obtiene el estado actual del procedimiento de ajuste angular. &
Status &
-- \\

GetErrorAxisID() &
Obtiene el identificador del eje responsable del error actual. &
UInt32 &
-- \\

GetAdjustingStatus() &
Obtiene el estado detallado del procedimiento de ajuste angular. &
AdjustingStatus &
-- \\

GetProfilePacketSumIndex(ProfileDataType profileDataType) &
Obtiene el número total de paquetes de datos disponibles para el tipo de perfil especificado. &
UInt32 &
-- \\

RequestProfileData(ProfileDataType profileDataType, UInt32 index) &
Solicita un paquete de datos correspondiente al perfil y al índice especificados. &
ProfileData &
-- \\

SetParameter(AngleAdjustmentParameter parameter) &
Configura los parámetros utilizados durante el procedimiento de ajuste angular. &
ErrorCode &
-- \\

SetParameter(AngleAdjustmentParameter parameter, RotationCenter rotationCenter) &
Configura los parámetros del procedimiento utilizando un centro de rotación definido. &
ErrorCode &
-- \\

%===========================================================
% Tipos anidados
%===========================================================
\SetRow{bg=gray9}
\SetCell[c=4]{c}
Tipos anidados \\

ErrorCode &
Enumeración que define los posibles códigos de error del procedimiento de ajuste angular. &
Enum &
-- \\

Status &
Enumeración que representa los estados posibles del procedimiento de ajuste angular. &
Enum &
-- \\

AdjustingStatus &
Enumeración que describe el estado interno del algoritmo de ajuste angular. &
Enum &
-- \\

ProfileDataType &
Enumeración que identifica los diferentes tipos de perfiles generados durante el procedimiento de ajuste angular. &
Enum &
-- \\

AngleAdjustmentParameter &
Clase que define los parámetros utilizados para configurar el procedimiento de ajuste angular. &
Class &
-- \\

ProfileData &
Clase que almacena los datos de perfil obtenidos durante el procedimiento de ajuste angular. &
Class &
-- \\

\end{longtblr}

\normalsize

%================================================================
%================================================================
%================================================================

% %===========================================================
% %===========================================================
% %                     Clase Profile
% %===========================================================
% %===========================================================

% \section{Clase Profile}



% %===========================================================
% %===========================================================
% %                     Clase Profile
% %===========================================================
% %===========================================================

% \section{Clase Profile}







\end{document}


